%======================================================================
% PRESENTACIÓN
%======================================================================
\chapter*{Presentación}
\addcontentsline{toc}{chapter}{Presentación}

\bigskip

\begin{cajaInfo}
El presente documento forma parte del ecosistema \textbf{Sandra}, una familia de sistemas de middleware empresarial desarrollados bajo principios de soberanía tecnológica.
\end{cajaInfo}

\bigskip

\textbf{Sandra Sentinel} representa la implementación de un principio fundamental en la ingeniería de sistemas modernos: la capacidad de verificar, auditar y recalcular de manera determinista los procesos críticos de negocio. En el contexto específico de la gestión de nóminas masivas en entornos jerárquicos, esta capacidad adquiere una relevancia especial dado el impacto directo en las finanzas del personal y en la responsabilidad fiscal de las organizaciones.

Este manual técnico ha sido diseñado para servir como referencia completa para ingenieros, arquitectos de sistemas, y administradores que interactúan con el sistema. El enfoque adoptado combina la rigurosidad matemática necesaria para comprender los algoritmos de cálculo con la praticidad requerida para la operación cotidiana del sistema.

\bigskip

\section*{Alcance del Manual}

Este documento aborda:

\begin{itemize}
    \item La arquitectura general del sistema y sus componentes.
    \item El modelo de datos y las entidades fundamentales.
    \item Los algoritmos de fusión mediante Hash Join.
    \item El motor de cálculo de nómina y sus fórmulas.
    \item El sistema de distribución de aportes con garantía de precisión.
    \item Los manifiestos de ejecución y su configuración.
    \item Detalles de implementación y consideraciones de despliegue.
\end{itemize}

\bigskip

\section*{Convenciones Utilizadas}

A lo largo de este documento se utilizan las siguientes convenciones:

\begin{itemize}
    \item \textbf{Fórmulas matemáticas:} Se presentan en notación matemática formal, siguiendo las convenciones de la literatura científica.
    \item \textbf{Código fuente:} Se presenta en tipografía monoespaciada con resaltado de sintaxis.
    \item \textbf{Diagramas:} Los diagramas de flujo y secuencia siguen la notación UML adaptada para sistemas de información.
    \item \textbf{Términos técnicos:} Los términos en inglés se presentan en cursiva la primera vez que aparecen, seguidos de su traducción al español.
\end{itemize}

\bigskip

\begin{cajaFormula}
Este símbolo indica fórmulas matemáticas importantes que definen el comportamiento del sistema.
\end{cajaFormula}

\vfill
\newpage
